\section{From distribution tests to finite graphs}\label{sec:graph-transfer}

We now prove \Cref{thm:base-graphs}.  The only probabilistic input to
this section is \Cref{thm:distribution-tests}, whose two parts were
proved in \Cref{prop:coloured-distribution,prop:asymmetric-distribution}.
The sample uses the thinned law \(\pi\).  Write
\begin{equation}\label{eq:transfer-scales}
 Q_N=2^{3000gN},\qquad \epsilon_N=2^{-100gN},\qquad
 m=2^{1000gN},\qquad \Omega=\supp\pi.
\end{equation}
All spaces in this section are finite.  In particular, \(\pi\) is
strictly positive on \(\Omega\), and
\(\log|\Omega|=O(N^2)\): specifying all frame entries uses
\(O(N^2)\) bits, and restricting their support by thinning only
decreases its size.  We regard the fixed construction parameters as
constants in this estimate.

\subsection{Entropy and terminal covers}

For probabilities \(p,q\) on a finite space with \(q>0\), define
\[
 \mathsf D(p\Vert q)=\sum_xp(x)\log\frac{p(x)}{q(x)},
 \qquad 0\log0=0.
\]
Logarithms in entropy expressions are natural logarithms.

This is the finite convex information-projection inequality underlying
the entropy method of \citet{csiszar1975}.

\begin{lemma}[Entropy increase after a deletion]\label{gt:entropy}
Let \(\mathcal P\) be a nonempty compact convex family of probabilities
and let \(p\) minimize \(\mathsf D(\cdot\Vert q)\) on \(\mathcal P\).
Then \(p\) is positive on the union of the feasible supports.  For
every \(p'\in\mathcal P\) supported on \(S\),
\[
 \mathsf D(p'\Vert q)-\mathsf D(p\Vert q)
 \ge \mathsf D(p'\Vert p)\ge-\log p(S).
\]
\end{lemma}
\begin{proof}
The minimum exists by continuity.  If \(p(x)=0<p'(x)\), the
right derivative of entropy along \((1-t)p+tp'\) is \(-\infty\)
at zero; all terms on the support of \(p\) have finite derivative.
This contradicts minimality.  Differentiating now on the common
feasible support gives
\[
 \sum_x(p'(x)-p(x))\log\frac{p(x)}{q(x)}\ge0.
\]
Subtracting the two entropies proves the first inequality.  If
\(p'\) is supported on \(S\), then
\[
 \mathsf D(p'\Vert p)
 =\mathsf D(p'\Vert p(\cdot\mid S))-\log p(S).
\]
Relative entropy is nonnegative: apply \(\log t\le t-1\) with
\(t=q(x)/p'(x)\) and sum on the positive support of \(p'\).
This proves the second inequality as well.
\end{proof}

We require a terminal cover for coloured units that accounts for
four different capacities.  The following elementary separation
argument supplies the finite linear programming duality needed here.

\begin{lemma}[Resource prices]\label{gt:prices}
Let \(\mathcal R\) be a nonempty finite collection of objects, each
having a nonnegative load vector \(v_r\in\R^d\), and let
\(c\in\R_{\ge0}^d\) be the vector of capacities.  If no probability
\((p_r)_{r\in\mathcal R}\) has \(\sum_rp_rv_r\le c\)
coordinatewise, there are nonnegative prices \(w\in\R^d\) with
\[
 w\cdot c<1,\qquad w\cdot v_r\ge1\quad(r\in\mathcal R).
\]
\end{lemma}
\begin{proof}
The compact convex sets
\(C=\operatorname{conv}\{v_r:r\in\mathcal R\}\) and
\(B=\prod_{j=1}^d[0,c_j]\) are disjoint.  Choose a closest pair
\(z\in C,b\in B\), and put \(u=z-b\ne0\).  For fixed \(z\),
the nearest point of \(B\) is its coordinatewise truncation, so
\(b_j=\min(z_j,c_j)\), \(u\ge0\), and \(u\cdot b=u\cdot c\).
For every \(v\in C\), comparison with \((1-t)z+tv\) at
\(t=0\) gives \(u\cdot(v-z)\ge0\).  Consequently
\[
 u\cdot v\ge u\cdot z
   =u\cdot c+\|u\|^2>u\cdot c\ge0.
\]
Take \(w=u/(u\cdot z)\).  This proves both assertions using
only minimization on compact subsets of a Euclidean space.
\end{proof}

\begin{lemma}[The coloured terminal cover]\label{gt:coloured-cut}
Fix a palette \([q]=\{1,\ldots,q\}\).  Suppose that a set
\(R\subseteq\Omega^2\times[q]\) supports no probability \(\sigma\)
whose two endpoint marginals are at most \(M\pi\), whose projection
onto \(\Omega^2\) is at most \(Q_N\pi^2\), and whose mass on
each colour is at most \(\eta_{\mathrm{col}}\).
There are sets \(S\subseteq\Omega\), \(E\subseteq\Omega^2\)
and \(C\subseteq[q]\) such that
\begin{equation}\label{eq:coloured-terminal-cover}
 \pi(S)<\frac8M,\qquad
 \pi^2(E)<\frac4{Q_N},\qquad
 |C|<\frac4{\eta_{\mathrm{col}}},
\end{equation}
and each \((x,y,c)\in R\) has \(x\in S\), \(y\in S\),
\((x,y)\in E\), or \(c\in C\).
\end{lemma}
\begin{proof}
The assertion is immediate for empty \(R\).  Otherwise introduce
separate resources for each left vertex, right vertex, ordered pair,
and colour, with capacities
\[
 M\pi(x),\quad M\pi(y),\quad Q_N\pi(x)\pi(y),\quad
 \eta_{\mathrm{col}},
\]
respectively.  An object \((x,y,c)\) uses one unit of each of its
four resources.  Apply \Cref{gt:prices}.  Write the resulting prices
as \(\ell_x,r_y,e_{xy},d_c\); their capacity-weighted sum is less
than one, and
\(\ell_x+r_y+e_{xy}+d_c\ge1\) on \(R\).
Take all resources of price at least \(1/4\).  They cover \(R\).
Let \(S\) be the union of the high-price left and right vertex
sets.  In fact
\[
 \pi(S)\le4\sum_x\pi(x)(\ell_x+r_x)<4/M<8/M.
\]
The corresponding pair and colour inequalities are
\(\pi^2(E)<4/Q_N\) and \(|C|<4/\eta_{\mathrm{col}}\).
This uses the capacity on the pair projection, so all colours of
one pair share a single pair resource.
\end{proof}

The next cover is a capacitated max-flow/min-cut argument
\citep{fordfulkerson1956}; we include the finite proof for the real
capacities used here.

\begin{lemma}[The uncoloured terminal cover]\label{gt:flow-cut}
If \(R\subseteq\Omega^2\) supports no probability with both
marginals at most \(L\pi\) and joint law at most \(Q_N\pi^2\),
then there are \(S\subseteq\Omega\) and \(E\subseteq\Omega^2\)
with
\begin{equation}\label{eq:uncoloured-terminal-cover}
 \pi(S)<1/L,\qquad \pi^2(E)<1/Q_N,
\end{equation}
such that each pair in \(R\) meets \(S\) or belongs to \(E\).
A set \(T\subseteq\Omega\) supports no probability at most
\(L\pi\) precisely when \(\pi(T)<1/L\).
\end{lemma}
\begin{proof}
Make a directed network from a source to a left copy of \(\Omega\),
then to a right copy, then to a sink.  The source and sink arcs
have capacities \(L\pi(x)\) and \(L\pi(y)\).  The middle arc
\(x\to y\), present for \((x,y)\in R\), has capacity
\(Q_N\pi(x)\pi(y)\).  A flow of value one is exactly a feasible
unit law, and a larger flow can be scaled to value one.

A maximum flow exists by compactness.  Its value is less than one.
In its residual network there is no source-to-sink path, since
augmenting by the smallest residual capacity on such a path would
increase the value.  Let \(Z\) be the vertices reachable from the
source.  Every original arc leaving \(Z\) is saturated; every
original arc entering \(Z\) carries zero flow.  Summing flow
conservation over \(Z\) therefore shows that its cut capacity is
the maximum flow value, hence is less than one.

Let \(A\) be the left types outside \(Z\), and \(B\) the right
types inside \(Z\).  The cut inequality reads
\[
 L\pi(A)+L\pi(B)
 +Q_N\pi^2\bigl(R\cap((\Omega\setminus A)
                         \times(\Omega\setminus B))\bigr)<1.
\]
Take \(S=A\cup B\) and take \(E\) to be the pair set in the
last summand.  This proves the cover and both bounds.  For the
point assertion, a law of total mass one on \(T\) bounded by
\(L\pi\) exists exactly when \(L\pi(T)\ge1\): necessity follows
by summing, and sufficiency follows by normalizing \(\pi|_T\).
\end{proof}

\subsection{Finite fingerprints}

Call two coloured units adjacent when their colours differ and
their four cross pairs are holes.  This is a finite simple graph:
a unit cannot conflict with itself, because a vertex has no hole
to itself.  For a fixed palette, consider its independent sets.
Also consider the bipartite graph with raw units on one side and
points of \(\Omega\) on the other, with adjacency meaning an
edge--point conflict.  In this graph the relevant objects are pairs
\((I,J)\) having no adjacency between \(I\) and \(J\).

The recording-and-replay construction follows the graph fingerprint
method of Kleitman and Winston \citep{kleitmanwinston1982}; see also
\citet{samotij2015}. Its potential here is the constrained entropy
minimum, whose increase makes the terminal count small.

\begin{lemma}[Fingerprint bound]\label{gt:fingerprints}
There are families of terminal sets, respectively terminal pairs
of sets, containing every such independent set, respectively every
such pair \((I,J)\), with the following properties.  Each coloured
terminal has an infeasible coloured region as in
\Cref{gt:coloured-cut}; in each bipartite terminal at least one of
the two feasible regions in \Cref{gt:flow-cut} is empty.
The number of terminals of either kind is
\begin{equation}\label{eq:transfer-terminal-count}
 \exp\bigl(O(N^3 2^{100gN})\bigr)=\exp(o(m)).
\end{equation}
The constants may depend on the fixed palette and on \(L\).
\end{lemma}
\begin{proof}
Start with the entire relevant graph as residual set.  On a
nonempty feasible region choose the entropy minimizer, breaking
all choices deterministically.  For coloured units the reference
probability is \(\pi^2\) times the uniform law on \([q]\).
The projected pair cap gives density at most \(qQ_N\) against
this reference.  The minimum entropy is therefore in
\([0,\log(qQ_N)]\).
By \Cref{thm:distribution-tests}(i), the average minimizer mass
of a residual neighbourhood is at least \(\epsilon_N\).
Select a vertex of maximal neighbourhood mass.  If it belongs
to the independent set being contained, record it and delete its
neighbours.  Otherwise delete only that vertex.  These operations
preserve containment of the target independent set.

For the bipartite problem, choose separate entropy minimizers on
the two residual sides, against \(\pi^2\) and \(\pi\).
Their sum lies in \([0,\log Q_N+\log L]\).
By \Cref{thm:distribution-tests}(ii), some vertex on either side
has opposite minimizer neighbourhood mass at least
\(\epsilon_N\).  Choose deterministically a vertex of maximal
such mass, including its side in the choice.  If it belongs to
the corresponding target set, record it and delete its neighbours
on the opposite side; otherwise delete only the chosen vertex.
The absence of edges between \(I\) and \(J\) preserves both
containments.

In either procedure every step removes a vertex: a recorded
vertex has a nonempty neighbourhood.  Stop when feasibility
fails (on either side in the bipartite case).  Deleting vertices
can only increase each minimum entropy.  If a recording leaves
the relevant feasible region nonempty, \Cref{gt:entropy} increases
its minimum by at least
\(-\log(1-\epsilon_N)\ge\epsilon_N\).
Consequently the number of recordings is at most
\begin{equation}\label{eq:transfer-fingerprint-length}
 1+\left\lceil
 \frac{\max\{\log(qQ_N),\log Q_N+\log L\}}
      {\epsilon_N}\right\rceil
 =O(N2^{100gN}).
\end{equation}
The extra one permits a last recording that destroys feasibility.

After recording a vertex, its residual neighbourhood is empty
forever.  Since a selected vertex always has positive neighbourhood
mass, it cannot be selected again.  The set of recorded vertices,
with side information in the bipartite case, thus determines the
entire procedure: replay its deterministic choices, performing a
recording exactly when the selected vertex is in that set, and
otherwise deleting it.  In particular it determines the terminal.
The total number of available vertices has logarithm \(O(N^2)\).
Bounding the number of fingerprints of length at most
\eqref{eq:transfer-fingerprint-length} by the corresponding power
of one plus that number proves \eqref{eq:transfer-terminal-count}.
\end{proof}

\subsection{One sample satisfies all three properties}

Choose the colour threshold \(\eta_{\mathrm{col}}>0\) from
\Cref{thm:distribution-tests}(i), and fix
\begin{equation}\label{eq:transfer-a-L}
 0<a<\min\{1/100,\eta_{\mathrm{col}}/800\},\qquad
 L>10^6/\beta.
\end{equation}
Here we may assume \(0<\beta\le1\), since increasing \(\beta\)
only weakens the required graph properties.  The order is
\(\eta_{\mathrm{col}},a\), then \(\beta,L\), followed by the
remaining construction parameters as in \Cref{sec:parameters}.
Let \(X_1,\ldots,X_m\) be independent samples from \(\pi\).
On positions \([m]\), put an edge between distinct positions
exactly when their types have no hole.  Repeated types are allowed.
By \Cref{lem:hole-triangle-free}, this graph has independence
number at most two for every sample.

First consider any matching whose edges have been partitioned
into classes of size at most \(am\).  Merge successive classes
until a merged class first has size at least \(am\), and repeat,
leaving a possible final smaller class.  Every merged class has
size at most \(2am\).  As a matching has at most \(m/2\) edges,
there are at most
\[
 q=\left\lceil\frac1{2a}\right\rceil+1
\]
merged classes.  Assign these the fixed palette \([q]\).
Different merged colours imply different original classes.

For every coloured terminal, fix a cover \((S,E,C)\) from
\Cref{gt:coloured-cut}.  Put \(k=\lceil m/100\rceil\).
The probability that at least \(k\) sampled positions lie in
\(S\) is at most
\[
 \binom mk(8/M)^k\le(800e/M)^k=\exp(-\Omega(m)).
\]
A fixed collection of \(k\) disjoint ordered position pairs
gives independent \(\pi^2\)-draws, and is entirely in \(E\)
with probability less than \((4/Q_N)^k\).  There are at most
\(m^{2k}\) such collections, so its union bound is
\begin{equation}\label{eq:transfer-pair-tail}
 m^{2k}(4/Q_N)^k=2^{-(1000gN-2)k}.
\end{equation}
Union over \eqref{eq:transfer-terminal-count}.  With probability
\(1-\exp(-\Omega(m))\), simultaneously every coloured terminal
has fewer than \(m/100\) exceptional positions and fewer than
\(m/100\) disjoint exceptional ordered pairs.  The strict bounds
use \(k-1<m/100\).

Suppose there were a matching of size at least \(m/20\) with no
conflict between original classes.  Orient its edges arbitrarily,
merge its classes as above, and take the set of resulting coloured
raw types.  This is an independent set in the coloured conflict
graph, so is contained in a terminal.  Every matching edge is
covered by an exceptional position, an exceptional ordered pair,
or an exceptional colour.  The first two covers each account for
fewer than \(m/100\) edges.  The last accounts for fewer than
\[
 \frac4{\eta_{\mathrm{col}}}\,2am
 <m/100
\]
edges.  Their sum is less than \(3m/100<m/20\), a contradiction.
This proves the first graph property.  Multiple sampled edges
may have the same raw type; the counting uses their disjoint
sample positions and does not assume distinct types.

For the bipartite terminals put \(k_\beta=\lceil\beta m/100\rceil\).
At a terminal with infeasible unit region, fix \((S,E)\) from
\Cref{gt:flow-cut}.  At a terminal with infeasible point region,
its entire point region has \(\pi\)-mass less than \(1/L\).
For every one of these vertex sets, the position tail is at most
\[
 \binom m{k_\beta}L^{-k_\beta}
 \le\left(\frac{100e}{\beta L}\right)^{k_\beta}
 =\exp(-\Omega_\beta(m)).
\]
The exceptional ordered-pair tail is at most
\[
 m^{2k_\beta}Q_N^{-k_\beta}=2^{-1000gNk_\beta}.
\]
Another union over \eqref{eq:transfer-terminal-count} shows that
all these sets simultaneously contain fewer than \(\beta m/100\)
sample positions, and all pair exceptions contain fewer than that
many disjoint ordered pairs, with probability tending to one.

If a matching and vertex set, each of size at least
\(\beta m/10\), had no edge--point conflict, their raw type sets
would give a pair \((I,J)\) contained in a bipartite terminal.
An infeasible point region cannot contain that vertex set.  An
infeasible unit region covers fewer than \(2\beta m/100\)
matching edges, also a contradiction.  We have proved the second
property even at the stronger threshold \(\beta m/10\).
There is no disjointness requirement between the original vertex
set and matching; a conflict itself ensures that its point differs
from both edge endpoints.

Finally let \(A,B\subseteq[m]\) each have size at least \(\beta m\).
A maximal matching in the induced graph on \(A\) leaves at most
two vertices unmatched: the unmatched vertices are independent,
and the whole graph has independence number at most two.
Its size is at least \((|A|-2)/2\ge\beta m/10\) for all
sufficiently large \(m\).  Apply the second property to this
matching and \(B\).  The resulting edge--point conflict supplies
a hole from \(A\) to \(B\), proving the third property.

All the required events hold together with probability tending
to one.  Their intersection therefore contains a finite sample
for every sufficiently large allowed \(n\); since
\(m=2^{1000gM_0n}\) tends to infinity, this proves
\Cref{thm:base-graphs}.
